%% paper66_informationstheorie_de.tex
%% Paper 66: Informationstheorie und Kodierungstheorie (Deutsch)
%% Autor: Michael Fuhrmann
%% Projekt: Specialist Mathematics Project, Build 145
%% Datum: 2026-03-12

\documentclass{amsart}

%% ---- Pakete ----
\usepackage[ngerman]{babel}
\usepackage{amsmath, amssymb, amsthm}
\usepackage{mathtools}
\usepackage{hyperref}
\usepackage{enumitem}
\usepackage{geometry}
\geometry{margin=1in}

%% ---- Theorem-Umgebungen (Deutsch) ----
\newtheorem{theorem}{Satz}[section]
\newtheorem{lemma}[theorem]{Lemma}
\newtheorem{korollar}[theorem]{Korollar}
\newtheorem{proposition}[theorem]{Proposition}
\theoremstyle{definition}
\newtheorem{definition}[theorem]{Definition}
\newtheorem{beispiel}[theorem]{Beispiel}
\newtheorem{bemerkung}[theorem]{Bemerkung}
\theoremstyle{remark}
\newtheorem{vermutung}[theorem]{Vermutung}
\newtheorem{offenefrage}[theorem]{Offene Frage}

%% ---- Operatoren ----
\DeclareMathOperator{\Ent}{H}
\DeclareMathOperator{\MI}{I}
\DeclareMathOperator{\KL}{D_{KL}}
\newcommand{\E}{\mathbb{E}}
\newcommand{\R}{\mathbb{R}}
\newcommand{\N}{\mathbb{N}}
\newcommand{\Z}{\mathbb{Z}}
\newcommand{\F}{\mathbb{F}}
\newcommand{\SNR}{\mathrm{SNR}}
\newcommand{\wt}{\mathrm{wt}}
\newcommand{\dist}{\mathrm{d}}

%% ---- Titelblock ----
\title{Informationstheorie und Kodierungstheorie:\\
Shannons mathematisches Fundament und seine Grundlagen}

\author{Michael Fuhrmann}
\address{Specialist Mathematics Project, Build 145}
\email{specialist-maths@project.local}
\date{2026-03-12}

%% ---- Abstract (VOR \maketitle!) ----
\begin{abstract}
Wir präsentieren eine strenge mathematische Darstellung der Informationstheorie
und Kodierungstheorie, basierend auf Claude Shannons grundlegendem Werk
\emph{A Mathematical Theory of Communication} (1948). Ausgehend von der
axiomatischen Definition der Shannon-Entropie
$H(X) = -\sum_i p_i \log_2 p_i$ entwickeln wir gegenseitige Information,
die Datenverarbeitungsungleichung und Shannons zwei fundamentale Sätze: den
Quellencodierungssatz (verlustlose Kompression) und den Kanalcodierungssatz
(Übertragung über verrauschte Kanäle). Wir beweisen die Optimalität des
Huffman-Codes, die Singleton-Schranke, die Hamming-Schranke sowie die
kapazitätserreichende Eigenschaft von Polarcodes (Ar\i{}kan 2009). Ferner
werden lineare Codes, Reed-Solomon-Codes, Turbo-Codes und LDPC-Codes behandelt.
Den Abschluss bildet die algorithmische Informationstheorie, in der die
Unberechenbarkeit der Kolmogorov-Komplexität $K(x)$ bewiesen wird. Alle
Ergebnisse sind mit ihrem Beweisstand ausgezeichnet: bewiesene Sätze versus
offene Vermutungen.
\end{abstract}

\maketitle

%% ============================================================
\section{Einleitung}
%% ============================================================

Im Jahr 1948 veröffentlichte Claude Elwood Shannon \emph{A Mathematical Theory
of Communication}~\cite{Shannon1948} -- ein Werk, das gleichzeitig die
Informationstheorie als mathematische Disziplin begründete und das
Ingenieurswissen der digitalen Kommunikation auf eine exakte Basis stellte.
Shannons zentraler Gedanke war, \emph{Information} als messbare Größe zu
quantifizieren, unabhängig vom semantischen Inhalt der Nachrichten, und die
fundamentalen Grenzen zuverlässiger Kommunikation über verrauschte Kanäle
zu charakterisieren.

Vor Shannon baute die Nachrichtentechnik auf Heuristiken und empirischem
Entwurf. Shannon bewies, dass jeder Kommunikationskanal eine präzise Kapazität
$C$ (gemessen in Bit pro Kanalnutzung) besitzt, und dass zuverlässige
Kommunikation bei jeder Rate $R < C$ erreichbar ist, während Raten $R > C$
grundsätzlich unmöglich sind. Dieses Resultat -- der Kanalcodierungssatz -- ist
eines der tiefgründigsten der angewandten Mathematik.

\subsection*{Überblick über dieses Paper}

Das vorliegende Paper behandelt:
\begin{enumerate}[label=(\arabic*)]
  \item Shannon-Entropie und ihre Eigenschaften (bewiesen);
  \item Gegenseitige Information und Datenverarbeitungsungleichung (bewiesen);
  \item Quellencodierungssatz und Optimalität des Huffman-Codes (bewiesen);
  \item Kanalkapazität und Shannon-Hartley-Formel für AWGN-Kanäle (bewiesen);
  \item Kanalcodierungssatz und starke Umkehrung (bewiesen);
  \item Fehlerkorrigierende Codes: Hamming- und Singleton-Schranke (bewiesen);
  \item Lineare Codes, Hamming-Codes, Reed-Solomon-Codes;
  \item Kapazitätserreichende Codes: Turbo, LDPC, Polar (Ar\i{}kan, bewiesen);
  \item Algorithmische Informationstheorie und Kolmogorov-Komplexität
        (teils bewiesen, teils offen).
\end{enumerate}

Alle Logarithmen sind zur Basis 2, sofern nicht anders angegeben; die
Informationseinheit ist das \emph{Bit}.

%% ============================================================
\section{Entropie}
%% ============================================================

\subsection{Definition und grundlegende Eigenschaften}

\begin{definition}[Shannon-Entropie]
Sei $X$ eine diskrete Zufallsvariable mit Werten in einem endlichen Alphabet
$\mathcal{X}$ und Wahrscheinlichkeitsverteilung $p(x) = \Pr[X = x]$.
Die \emph{Shannon-Entropie} von $X$ ist
\begin{equation}\label{eq:entropie}
  H(X) \;=\; -\sum_{x \in \mathcal{X}} p(x)\log_2 p(x),
\end{equation}
wobei wir die Konvention $0 \log_2 0 = 0$ verwenden (durch Stetigkeit gerechtfertigt).
Die Einheit ist \emph{Bit}.
\end{definition}

\begin{bemerkung}
Äquivalent gilt $H(X) = \E[-\log_2 p(X)]$, d.\,h.\ die erwartete
\emph{Überraschung} (Selbstinformation). Der \emph{Informationsgehalt} (Surprisal)
eines Ereignisses $x$ ist $I(x) = -\log_2 p(x)$.
\end{bemerkung}

\begin{beispiel}
Eine faire Münze ($p = 1/2$) hat $H = 1$ Bit. Eine verzerrte Münze mit $p = 0{,}9$
hat $H = -0{,}9 \log_2 0{,}9 - 0{,}1 \log_2 0{,}1 \approx 0{,}469$ Bit.
\end{beispiel}

\begin{theorem}[Nicht-Negativität]\label{satz:nichtneg}
$H(X) \geq 0$, mit Gleichheit genau dann, wenn $X$ deterministisch ist
(d.\,h.\ $p(x^*)=1$ für ein bestimmtes $x^*$).
\end{theorem}
\begin{proof}
Da $0 \leq p(x) \leq 1$, gilt $-\log_2 p(x) \geq 0$ für alle $x$ mit
$p(x)>0$. Daher $H(X) \geq 0$. Gleichheit gilt genau dann, wenn
$-p(x)\log_2 p(x)=0$ für alle $x$, d.\,h.\ $p(x) \in \{0,1\}$ für alle $x$.
\end{proof}

\begin{theorem}[Maximale Entropie bei Gleichverteilung]\label{satz:maxent}
Unter allen Verteilungen auf einem Alphabet $\mathcal{X}$ der Größe $n$ wird
die Entropie durch die Gleichverteilung $p(x) = 1/n$ maximiert, mit
$H(X) = \log_2 n$.
\end{theorem}
\begin{proof}
Nach der Log-Summen-Ungleichung (oder Jensen angewandt auf die konkave Funktion
$-t\log_2 t$) gilt für jede Verteilung $p$ auf $\mathcal{X}$:
\[
  H(X) = -\sum_x p(x)\log_2 p(x) \leq \log_2 n,
\]
mit Gleichheit genau bei $p(x) = 1/n$ für alle $x$. Konkret:
\[
  \log_2 n - H(X) = \sum_x p(x)\log_2\!\frac{p(x)}{1/n} = \KL(p \,\|\, u) \geq 0,
\]
wobei $u$ die Gleichverteilung ist und $\KL \geq 0$ nach der Gibbs-Ungleichung
(Nicht-Negativität der KL-Divergenz).
\end{proof}

\subsection{Verbundentropie und bedingte Entropie}

\begin{definition}
Für gemeinsam verteilte Zufallsvariablen $X, Y$ mit Verbundverteilung $p(x,y)$:
\begin{align}
  H(X,Y) &= -\sum_{x,y} p(x,y)\log_2 p(x,y), \\
  H(X \mid Y) &= -\sum_{x,y} p(x,y)\log_2 p(x \mid y) = H(X,Y) - H(Y).
\end{align}
\end{definition}

\begin{theorem}[Subadditivität der Entropie]\label{satz:subadditivitaet}
$H(X,Y) \leq H(X) + H(Y)$, mit Gleichheit genau dann, wenn $X$ und $Y$
unabhängig sind.
\end{theorem}
\begin{proof}
$H(X) + H(Y) - H(X,Y) = \MI(X;Y) \geq 0$
(gegenseitige Information ist nicht-negativ, bewiesen als Satz~\ref{satz:mi-nichtneg}).
\end{proof}

\begin{theorem}[Konditionierung verringert Entropie]\label{satz:kondred}
$H(X \mid Y) \leq H(X)$, mit Gleichheit genau dann, wenn $X \perp Y$.
\end{theorem}
\begin{proof}
$H(X) - H(X\mid Y) = \MI(X;Y) \geq 0$.
\end{proof}

%% ============================================================
\section{Gegenseitige Information}
%% ============================================================

\subsection{Definition}

\begin{definition}[Gegenseitige Information]
Die \emph{gegenseitige Information} (mutual information) zweier diskreter
Zufallsvariablen $X$ und $Y$ ist
\begin{equation}\label{eq:mi}
  \MI(X;Y) \;=\; H(X) + H(Y) - H(X,Y)
  \;=\; \sum_{x,y} p(x,y)\log_2\frac{p(x,y)}{p(x)p(y)}.
\end{equation}
Äquivalent: $\MI(X;Y) = \KL(p_{XY} \,\|\, p_X \otimes p_Y)$.
\end{definition}

\begin{theorem}[Nicht-Negativität der gegenseitigen Information]\label{satz:mi-nichtneg}
$\MI(X;Y) \geq 0$, mit Gleichheit genau dann, wenn $X$ und $Y$ unabhängig sind.
\end{theorem}
\begin{proof}
$\MI(X;Y) = \KL(p_{XY}\|p_Xp_Y)$. Nach der Gibbs-Ungleichung (Jensen auf
$-\log$):
\[
  \KL(p \| q) = \sum_x p(x)\log_2\frac{p(x)}{q(x)} \geq 0,
\]
mit Gleichheit genau dann, wenn $p = q$ überall. Daher $\MI(X;Y) \geq 0$, mit
Gleichheit genau dann, wenn $p(x,y) = p(x)p(y)$ für alle $x,y$.
\end{proof}

\subsection{Datenverarbeitungsungleichung}

\begin{theorem}[Datenverarbeitungsungleichung]\label{satz:dpu}
Bildet $X \to Y \to Z$ eine Markov-Kette (d.\,h.\ $Z$ ist bedingt unabhängig
von $X$ gegeben $Y$), so gilt
\[
  \MI(X;Z) \leq \MI(X;Y).
\]
\end{theorem}
\begin{proof}
Durch die Kettenregel der gegenseitigen Information:
\begin{align*}
  \MI(X;Y,Z) &= \MI(X;Y) + \MI(X;Z\mid Y), \\
  \MI(X;Y,Z) &= \MI(X;Z) + \MI(X;Y\mid Z).
\end{align*}
Da $X \to Y \to Z$ eine Markov-Kette ist, gilt $\MI(X;Z \mid Y) = 0$. Also
$\MI(X;Y) = \MI(X;Z) + \MI(X;Y \mid Z)$.
Da $\MI(X;Y\mid Z) \geq 0$, folgt $\MI(X;Z) \leq \MI(X;Y)$.
\end{proof}

\begin{bemerkung}
Die Datenverarbeitungsungleichung besagt, dass keine Nachbearbeitung von $Y$
die Information über $X$ erhöhen kann. Sie ist grundlegend im maschinellen Lernen
(keine deterministische Funktion der Merkmale kann die Information in den
Merkmalen übersteigen) und in der Netzwerkinformationstheorie.
\end{bemerkung}

%% ============================================================
\section{Quellencodierungssatz}
%% ============================================================

\subsection{Verlustlose Kompression}

\begin{definition}[Präfixfreier Code]
Ein \emph{präfixfreier} (instantaner) Code für $\mathcal{X}$ weist jedem
$x \in \mathcal{X}$ ein binäres Codewort $C(x) \in \{0,1\}^*$ zu, sodass kein
Codewort Präfix eines anderen ist. Die \emph{mittlere Codelänge} ist
$\bar{L} = \sum_x p(x)\,|C(x)|$.
\end{definition}

\begin{theorem}[Quellencodierungssatz (Shannon 1948)]\label{satz:quellencodierung}
Für jede diskrete Zufallsvariable $X$ mit Entropie $H(X)$ existiert ein
präfixfreier Code mit mittlerer Codelänge $\bar{L}$, sodass
\[
  H(X) \;\leq\; \bar{L} \;<\; H(X) + 1.
\]
Kein präfixfreier Code kann $\bar{L} < H(X)$ erreichen.
\end{theorem}
\begin{proof}
\textbf{Untere Schranke.} Nach der Kraft-Ungleichung erfüllt jeder präfixfreie
Code $\sum_x 2^{-|C(x)|} \leq 1$. Sei $l_x = |C(x)|$. Dann:
\begin{align*}
  \bar{L} - H(X) &= \sum_x p(x)(l_x + \log_2 p(x)) \\
  &= \sum_x p(x)\log_2\frac{p(x)}{2^{-l_x}}
   = \KL\!\left(p \,\Big\|\, \{2^{-l_x}\}\right) \geq 0.
\end{align*}

\textbf{Obere Schranke.} Wähle $l_x = \lceil -\log_2 p(x)\rceil$. Dann gilt
$l_x < -\log_2 p(x) + 1$, also $\sum_x 2^{-l_x} \leq \sum_x p(x) = 1$
(Kraft erfüllt). Die mittlere Codelänge erfüllt
\[
  \bar{L} = \sum_x p(x)l_x < \sum_x p(x)(-\log_2 p(x) + 1) = H(X) + 1. \qedhere
\]
\end{proof}

\subsection{Huffman-Codes}

\begin{definition}[Huffman-Algorithmus (1952)]
Gegeben Symbole mit Wahrscheinlichkeiten $p_1 \geq p_2 \geq \cdots \geq p_n$:
\begin{enumerate}
  \item Behandle jedes Symbol als Blattknoten mit Gewicht $p_i$.
  \item Verschmelze wiederholt die zwei Knoten mit kleinstem Gewicht zu einem
        neuen Elternknoten (Gewicht = Summe), beschrifte die Kinder mit 0 und 1.
  \item Der entstehende Binärbaum definiert einen präfixfreien Code.
\end{enumerate}
\end{definition}

\begin{theorem}[Optimalität des Huffman-Codes]\label{satz:huffman}
Unter allen präfixfreien Codes minimiert der Huffman-Code die mittlere
Codewortlänge $\bar{L}$.
\end{theorem}
\begin{proof}[Beweisschema]
Durch vollständige Induktion über $n$. Induktionsanfang $n=2$: beide Symbole
erhalten 1-Bit-Codes; optimal. Induktionsschritt: Sei $C^*$ ein beliebiger
optimaler Code. Unter allen Symbolen mit minimaler Wahrscheinlichkeit haben
mindestens zwei dieselbe maximale Codewortlänge (andernfalls ließe sich das
längere Wort verkürzen und $\bar{L}$ reduzieren). O.\,B.\,d.\,A.\ bilden
die zwei unwahrscheinlichsten Symbole $x_{n-1}, x_n$ in $C^*$ ein
Geschwisterpaar. Zusammenfassen zu einem Meta-Symbol mit Wahrscheinlichkeit
$p_{n-1}+p_n$ ergibt einen $(n-1)$-Symbol-Code $C'$ mit
$\bar{L}(C^*) = \bar{L}(C') + (p_{n-1}+p_n)$. Nach Induktionsvoraussetzung
minimiert Huffmans gieriges Zusammenfassen $\bar{L}(C')$, also auch
$\bar{L}(C^*)$.
\end{proof}

%% ============================================================
\section{Kanalkapazität}
%% ============================================================

\subsection{Diskrete gedächtnislose Kanäle}

\begin{definition}[Diskreter gedächtnisloser Kanal]
Ein \emph{diskreter gedächtnisloser Kanal} (DMC) ist durch das Eingangsalphabet
$\mathcal{X}$, das Ausgangsalphabet $\mathcal{Y}$ und Übergangswahrscheinlichkeiten
$p(y \mid x)$ gegeben. Gedächtnislos bedeutet: $p(y^n \mid x^n) =
\prod_{i=1}^n p(y_i \mid x_i)$.
\end{definition}

\begin{definition}[Kanalkapazität]
Die \emph{Kapazität} eines DMC ist
\begin{equation}\label{eq:kapazitaet}
  C \;=\; \max_{p(x)} \MI(X;Y),
\end{equation}
wobei das Maximum über alle Eingangsverteilungen $p(x)$ genommen wird.
\end{definition}

\begin{beispiel}[Binärer symmetrischer Kanal]
Ein BSC mit Übergangswahrscheinlichkeit $\varepsilon \in [0,1/2)$ hat
$C_{\mathrm{BSC}} = 1 - H_b(\varepsilon)$, wobei $H_b(\varepsilon) =
-\varepsilon\log_2\varepsilon - (1-\varepsilon)\log_2(1-\varepsilon)$
die binäre Entropiefunktion ist.
\end{beispiel}

\subsection{AWGN-Kanal und Shannon-Hartley-Theorem}

\begin{definition}[AWGN-Kanal]
Das \emph{additive weiße Gauß'sche Rauschen} (AWGN) Kanalmodell ist
$Y = X + Z$, wobei $Z \sim \mathcal{N}(0, N_0/2)$ unabhängig vom Eingang $X$
ist, der einer Leistungsschranke $\E[X^2] \leq P$ genügt.
\end{definition}

\begin{theorem}[Shannon-Hartley-Theorem]\label{satz:shannon-hartley}
Die Kapazität des AWGN-Kanals mit Signalleistung $P$ und Rauschleistung $N$ ist
\begin{equation}\label{eq:awgn}
  C \;=\; \tfrac{1}{2}\log_2\!\left(1 + \tfrac{P}{N}\right)
  \quad [\text{Bit pro Kanalnutzung}].
\end{equation}
Für eine Bandbreite $W$ (Hz) und Rauschleistungsdichte $N_0/2$ ergibt sich
$C = W \log_2(1 + P/(N_0 W))$ Bit pro Sekunde.
\end{theorem}
\begin{proof}[Beweisschema]
Die Entropie-Leistungs-Ungleichung besagt, dass die Gauß-Verteilung die
differentielle Entropie $h(Y)$ unter Leistungsschranke maximiert. Mit
$h(Z) = \frac{1}{2}\log_2(2\pi e N)$ und $h(Y) \leq \frac{1}{2}\log_2(2\pi
e(P+N))$ folgt $\MI(X;Y) = h(Y) - h(Z) \leq \frac{1}{2}\log_2(1+P/N)$,
angenommen von $X \sim \mathcal{N}(0,P)$.
\end{proof}

%% ============================================================
\section{Kanalcodierungssatz}
%% ============================================================

\begin{theorem}[Kanalcodierungssatz (Shannon 1948)]\label{satz:kanalcodierung}
Für einen DMC mit Kapazität $C$ gilt:
\begin{enumerate}[label=(\alph*)]
  \item \textbf{Erreichbarkeit.} Für jede Rate $R < C$ und jedes $\varepsilon > 0$
        existieren Blockcodes der Länge $n$ (für genügend großes $n$) mit Rate $R$
        und maximaler Fehlerwahrscheinlichkeit $P_e^{(n)} < \varepsilon$.
  \item \textbf{Umkehrung.} Falls $R > C$, dann gilt für jede Folge von Codes
        $P_e^{(n)} \to 1$ für $n \to \infty$.
\end{enumerate}
\end{theorem}
\begin{proof}[Beweisschema (Zufällige Codierung)]
\textbf{Erreichbarkeit.} Erzeuge ein Codebuch $\mathcal{C}$ mit $2^{nR}$
Codewörtern $\mathbf{x}(m)$, $m \in [2^{nR}]$, jedes i.\,i.\,d.\ aus der
kapazitätserreichenden Eingangsverteilung $p^*(x)$. Der Decoder verwendet
\emph{typische Mengendekodierung}: erkläre $\hat{m}$, falls
$(\mathbf{x}(\hat{m}), \mathbf{y})$ gemeinsam typisch und eindeutig ist.
Nach dem asymptotischen Gleichverteilungssatz (AEP) ist die Wahrscheinlichkeit,
dass das korrekte Codewort typisch ist, $\geq 1-\varepsilon$ für großes $n$.
Die Wahrscheinlichkeit, dass ein falsches Codewort $m' \neq m$ gemeinsam typisch
mit $\mathbf{y}$ ist, beträgt $\leq 2^{-n(I(X;Y)-\varepsilon)}$. Per Union Bound
über $2^{nR}$ falsche Codewörter:
\[
  P_e \leq \varepsilon + 2^{nR} \cdot 2^{-n(C-\varepsilon)} \to 0 \quad \text{falls } R < C.
\]

\textbf{Umkehrung.} Durch Fano's Ungleichung: bei kleinem $P_e$ gilt
$H(M \mid Y^n) \leq 1 + P_e \cdot nR$. Außerdem $\MI(M; Y^n) \leq nC$ (Kanalschranke).
Da $nR = H(M) = \MI(M;Y^n) + H(M\mid Y^n)$, folgt $R \leq C + O(P_e)$,
was $P_e \not\to 0$ erzwingt, wenn $R > C$.
\end{proof}

\begin{theorem}[Starke Umkehrung]\label{satz:starke-umkehrung}
Für einen DMC mit Kapazität $C$ gilt: Falls $R > C$, so gilt
$P_e^{(n)} \to 1$ exponentiell schnell.
\end{theorem}
\begin{proof}[Beweisschema]
Mit der Methode der Typen oder Kugelpackungsexponenten erfüllt die
Fehlerwahrscheinlichkeit $P_e^{(n)} \geq 1 - 2^{-n(E(R) - o(1))}$ für einen
positiven Exponenten $E(R)>0$ bei $R>C$, was $P_e^{(n)} \to 1$ zeigt.
\end{proof}

%% ============================================================
\section{Fehlerkorrigierende Codes}
%% ============================================================

\subsection{Hamming-Abstand und Grunddefinitionen}

\begin{definition}[Hamming-Abstand]
Für $\mathbf{u}, \mathbf{v} \in \{0,1\}^n$ ist der \emph{Hamming-Abstand}
\[
  \dist_H(\mathbf{u}, \mathbf{v}) = |\{i : u_i \neq v_i\}|.
\]
Das \emph{Hamming-Gewicht} von $\mathbf{v}$ ist $\wt(\mathbf{v}) =
\dist_H(\mathbf{v}, \mathbf{0})$.
\end{definition}

\begin{definition}[$[n,k,d]$-Code]
Ein \emph{binärer Blockcode} $\mathcal{C} \subseteq \{0,1\}^n$ mit
$|\mathcal{C}| = 2^k$ und Minimaldistanz $d = \min_{\mathbf{c} \neq \mathbf{c}'}
\dist_H(\mathbf{c}, \mathbf{c}')$ ist ein \emph{$[n,k,d]$-Code}. Er kann:
\begin{itemize}
  \item bis zu $d-1$ Fehler erkennen,
  \item bis zu $\lfloor (d-1)/2 \rfloor$ Fehler korrigieren.
\end{itemize}
Die \emph{Rate} ist $R = k/n$. Die \emph{Redundanz} beträgt $n - k$ Bit.
\end{definition}

\subsection{Singleton-Schranke}

\begin{theorem}[Singleton-Schranke]\label{satz:singleton}
Jeder $[n,k,d]$-Code erfüllt
\[
  k \;\leq\; n - d + 1.
\]
\end{theorem}
\begin{proof}
Betrachte die $2^k$ Codewörter. Projiziere auf die ersten $k' = n-d+1$
Koordinaten; die entstehenden Vektoren müssen alle verschieden sein (andernfalls
stimmen zwei Codewörter in allen $k'$ Positionen überein, und da sie
Abstand $\geq d$ haben, müssten sie sich in mindestens $d$ der verbleibenden
$n-k'=d-1$ Positionen unterscheiden, was unmöglich ist). Daher
$2^k \leq 2^{k'} = 2^{n-d+1}$.
\end{proof}

\begin{definition}[MDS-Codes]
Codes, die $k = n - d + 1$ erreichen, heißen \emph{Maximum Distance Separable}
(MDS). Reed-Solomon-Codes sind MDS.
\end{definition}

\subsection{Hamming-Schranke (Kugelpackungsschranke)}

\begin{theorem}[Hamming-Schranke]\label{satz:hamming-schranke}
Ein $[n,k,d]$-Code mit $d \geq 2t+1$ (der $t$ Fehler korrigiert) erfüllt
\[
  2^k \cdot \sum_{i=0}^{t}\binom{n}{i} \;\leq\; 2^n.
\]
\end{theorem}
\begin{proof}
Jedes Codewort $\mathbf{c}$ hat eine \emph{Hamming-Kugel} vom Radius $t$:
$B(\mathbf{c}, t) = \{\mathbf{v} : \dist_H(\mathbf{v},\mathbf{c}) \leq t\}$
mit $|B(\mathbf{c},t)| = \sum_{i=0}^t \binom{n}{i}$. Da $d \geq 2t+1$, sind
diese Kugeln disjunkt und alle Teilmengen von $\{0,1\}^n$. Abzählen liefert
die Schranke.
\end{proof}

\begin{definition}[Perfekter Code]
Ein Code, der die Hamming-Schranke mit Gleichheit erfüllt, heißt
\emph{perfekter Code}. Hamming-Codes (siehe Abschnitt~\ref{sec:linear-de})
sind perfekte Codes.
\end{definition}

%% ============================================================
\section{Lineare Codes}\label{sec:linear-de}
%% ============================================================

\subsection{Generatormatrix und Prüfmatrix}

\begin{definition}[Linearer Code]
Ein \emph{linearer $[n,k]$-Code} über $\mathbb{F}_2$ ist ein $k$-dimensionaler
Unterraum von $\mathbb{F}_2^n$. Er wird beschrieben durch:
\begin{itemize}
  \item Eine \emph{Generatormatrix} $G \in \mathbb{F}_2^{k \times n}$: Codewörter
        sind $\mathbf{c} = \mathbf{m} G$ für Nachrichten $\mathbf{m} \in \mathbb{F}_2^k$.
  \item Eine \emph{Prüfmatrix} (Parity-Check Matrix) $H \in \mathbb{F}_2^{(n-k)
        \times n}$ mit $H G^T = 0$: $\mathbf{c}$ ist ein Codewort genau dann, wenn
        $H \mathbf{c}^T = \mathbf{0}$.
\end{itemize}
\end{definition}

\begin{definition}[Syndrom]
Für ein empfangenes Wort $\mathbf{r}$ ist das \emph{Syndrom}
$\mathbf{s} = H\mathbf{r}^T$. Falls $\mathbf{s} = \mathbf{0}$, wurde kein
(erkennbarer) Fehler übertragen.
\end{definition}

\subsection{Hamming-Codes}

\begin{definition}[Hamming-Code $[2^r-1, 2^r-r-1, 3]$]
Für $r \geq 2$ hat der \emph{Hamming-Code} eine Prüfmatrix $H$, deren Spalten
alle $2^r - 1$ von Null verschiedenen Binärvektoren der Länge $r$ sind. Parameter:
$n = 2^r - 1$, $k = 2^r - r - 1$, $d = 3$.
\end{definition}

\begin{proposition}[Hamming-Codes sind perfekt]
Hamming-Codes erfüllen die Hamming-Schranke mit Gleichheit und sind damit
perfekte 1-Fehler-korrigierende Codes.
\end{proposition}
\begin{proof}
Mit $t=1$: $2^k \cdot (1 + n) = 2^{2^r-r-1} \cdot 2^r = 2^{2^r-1} = 2^n$.
\end{proof}

\subsection{Reed-Solomon-Codes}

\begin{definition}[Reed-Solomon-Code]
Sei $q = p^m$ eine Primzahlpotenz und $\alpha_1,\ldots,\alpha_n$ seien $n$
verschiedene Elemente von $\mathbb{F}_q$. Der \emph{Reed-Solomon-Code} der
Dimension $k$ ist
\[
  \mathcal{C}_{RS} = \bigl\{(f(\alpha_1),\ldots,f(\alpha_n)) :
  f \in \mathbb{F}_q[x],\; \deg f \leq k-1\bigr\}.
\]
Es handelt sich um einen $[n, k, n-k+1]$-Code über $\mathbb{F}_q$, also MDS.
\end{definition}

\begin{bemerkung}
Reed-Solomon-Codes über $\mathbb{F}_{2^8}$ werden für die Fehlerkorrektur auf
CDs/DVDs und in RAID-6-Speichersystemen eingesetzt. Effiziente Dekodierung bis
zu $\lfloor(n-k)/2\rfloor$ Fehlern ist durch den Berlekamp-Massey- oder
euklidischen Algorithmus möglich.
\end{bemerkung}

%% ============================================================
\section{Kapazitätserreichende Codes}
%% ============================================================

\subsection{Turbo-Codes}

\begin{definition}[Turbo-Code (Berrou--Glavieux--Thitimajshima, 1993)]
Ein \emph{Turbo-Code} ist eine Parallelverkettung zweier rekursiver systematischer
Faltungscodierer (RSC), getrennt durch einen Pseudo-Zufalls-Interleaver. Die
Dekodierung erfolgt iterativ durch \emph{Belief Propagation} (Turbo-Dekodierung)
zwischen den zwei Komponentendecodern.
\end{definition}

\begin{bemerkung}
Turbo-Codes waren die ersten praxistauglichen Codes, die nahe an die
Shannon-Kapazität für AWGN-Kanäle herankamen. Die kapazitätserreichende
Eigenschaft im Grenzwert großer Blocklängen folgt aus der Dichteentwicklungs-Analyse.
\end{bemerkung}

\subsection{LDPC-Codes}

\begin{definition}[LDPC-Code (Gallager, 1960/1963)]
Ein \emph{Low-Density Parity-Check (LDPC)}-Code ist ein linearer Code, dessen
Prüfmatrix $H$ dünn besetzt ist (die meisten Einträge sind Null). Gallagers
ursprüngliche reguläre LDPC-Codes haben konstante Spaltengewicht $d_v$ und
Zeilengewicht $d_c$.
\end{definition}

\begin{bemerkung}
LDPC-Codes, dekodiert durch Belief Propagation auf dem Tanner-Graphen, können
Raten beliebig nahe an der Shannon-Kapazität des AWGN-Kanals erreichen.
Dies wurde von Richardson und Urbanke (2001) durch Dichteentwicklung gezeigt.
\end{bemerkung}

\subsection{Polarcodes und Ar\i{}kans Satz}

\begin{definition}[Polarcode (Ar\i{}kan, 2009)]
Sei $W$ ein binärer gedächtnisloser Kanal mit Kapazität $C(W)$. Polarcodes
nutzen das Phänomen der \emph{Kanalpolarisierung}: Durch Anwenden von $n = 2^\ell$
Kopien von $W$ unter der Polarisierungstransformation $G_n = B_n F^{\otimes \ell}$
(mit $F = \begin{pmatrix}1&0\\1&1\end{pmatrix}$ und $B_n$ als Bitreihenfolgeumkehrung)
entstehen $n$ synthetische Kanäle $W_n^{(i)}$, $i = 1, \ldots, n$, die sich zu
entweder perfekten ($I(W_n^{(i)}) \to 1$) oder nutzlosen ($I(W_n^{(i)}) \to 0$)
Kanälen polarisieren. Informationsbits werden nur auf den ``guten'' Kanälen gesendet;
die restlichen werden eingefroren.
\end{definition}

\begin{theorem}[Kapazitätserreichende Eigenschaft von Polarcodes (Ar\i{}kan 2009)]
\label{satz:polar}
Für jeden binären gedächtnislosen symmetrischen Kanal $W$ mit Kapazität $C(W)$
und jede Rate $R < C(W)$ erreichen Polarcodes der Blocklänge $n = 2^\ell$ eine
Blockfehlerwahrscheinlichkeit
\[
  P_e \;\leq\; 2^{-n^\beta}
\]
für jedes $\beta < 1/2$ und hinreichend großes $n$, bei Codier- und
Dekodierkomplexität $O(n \log n)$.
\end{theorem}
\begin{proof}[Beweisschema]
Ar\i{}kans Kernidee ist das Polarisierungslemma: für $I_n = \frac{1}{n}\sum_i
I(W_n^{(i)})$ gilt $I_n = C(W)$ für alle $n$ (Kapazitätserhaltung). Durch den
Martingal-Konvergenzsatz konvergiert der Anteil der Kanäle mit $I(W_n^{(i)})
> 1-\delta$ gegen $C(W)$, und der Anteil mit $I(W_n^{(i)}) < \delta$ gegen
$1-C(W)$. Die Fehlerwahrscheinlichkeit auf jedem guten Kanal fällt wie
$2^{-2^{\ell\beta}}$ durch die rekursive Struktur der sukzessiven Auslöschungsdekodierung,
woraus die behauptete Schranke folgt. Details in \cite{Arikan2009}.
\end{proof}

\begin{bemerkung}
Polarcodes sind die ersten provably kapazitätserreichenden Codes mit effizienter
($O(n\log n)$) Kodierung und Dekodierung. Sie sind im 5G-NR-Standard für die
Steuerkanal-Codierung eingesetzt.
\end{bemerkung}

%% ============================================================
\section{Algorithmische Informationstheorie}
%% ============================================================

\subsection{Kolmogorov-Komplexität}

\begin{definition}[Kolmogorov-Komplexität (Solomonoff 1960, Kolmogorov 1965, Chaitin 1969)]
Sei $\mathcal{U}$ eine universelle Turing-Maschine. Die \emph{Kolmogorov-Komplexität}
einer endlichen Binärzeichenkette $x$ ist
\[
  K(x) \;=\; \min\bigl\{|\pi| : \mathcal{U}(\pi) = x\bigr\},
\]
die Länge des kürzesten Programms $\pi$, das $\mathcal{U}$ veranlasst, $x$
auszugeben.
\end{definition}

\begin{theorem}[Invarianz]\label{satz:invarianz}
Bis auf eine nur von der Wahl der universellen Turing-Maschine abhängige additive
Konstante ist $K(x)$ wohldefiniert (unabhängig von $\mathcal{U}$).
\end{theorem}
\begin{proof}
Jede universelle Turing-Maschine $\mathcal{U}_1$ kann $\mathcal{U}_2$ mit einem
festen Präfix (Interpreter) der Länge $c_{12}$ simulieren. Daher
$K_1(x) \leq K_2(x) + c_{12}$, und symmetrisch $K_2(x) \leq K_1(x) + c_{21}$.
\end{proof}

\begin{theorem}[Unberechenbarkeit der Kolmogorov-Komplexität]\label{satz:kolmogorov-unberechenbar}
Die Funktion $x \mapsto K(x)$ ist nicht berechenbar.
\end{theorem}
\begin{proof}
Angenommen, $K$ wäre berechenbar. Ordne alle Binärzeichenketten $x_1, x_2, \ldots$
lexikographisch und sei $n$ eine große ganze Zahl. Definiere die Zeichenkette
$s$ = ``die erste Zeichenkette $x$ mit $K(x) > n$''. Diese Zeichenkette $s$ ist
wohldefiniert (da $K$ unbeschränkt wächst) und berechenbar aus $n$. Das Programm,
das $s$ berechnet, hat Länge $O(\log n)$, also $K(s) = O(\log n) \ll n$ für
großes $n$. Andererseits gilt per Definition $K(s) > n$ -- ein Widerspruch.
\end{proof}

\subsection{Chaitins $\Omega$ und algorithmische Zufälligkeit}

\begin{definition}[Chaitins $\Omega$]
Für eine präfixfreie universelle Turing-Maschine ist die \emph{Haltwahrscheinlichkeit}
\[
  \Omega \;=\; \sum_{\mathcal{U}(\pi)\text{ hält}} 2^{-|\pi|}
\]
eine reelle Zahl in $(0,1)$, die \emph{algorithmisch zufällig} ist: ihre binäre
Entwicklung ist inkomprimierbar in dem Sinne, dass
$K(\Omega \upharpoonright n) \geq n - O(1)$.
\end{definition}

\begin{offenefrage}[Chaitins $\Omega$]
Keine explizite geschlossene Darstellung von $\Omega$ ist bekannt. Kein Algorithmus
kann mehr als endlich viele Bits von $\Omega$ berechnen; tatsächlich kodiert jedes
Bit von $\Omega$ eine unabhängige ``orakulare'' Aussage über das Halteproblem.
Ob weitere strukturelle Eigenschaften von $\Omega$ bestimmbar sind, bleibt offen.
\end{offenefrage}

\begin{theorem}[AIT und Shannon-Entropie (Brudno 1983)]
Für eine ergodische stationäre Quelle $\mu$ gilt für fast jede Folge $x$:
\[
  \lim_{n\to\infty} \frac{K(x_{1:n})}{n} = H(\mu),
\]
wobei $H(\mu)$ die Entropierate der Quelle ist.
\end{theorem}
\begin{proof}
Dies folgt aus Brudnos Theorem, das die Kolmogorov-Komplexitätsrate einzelner
Folgen mit der maßtheoretischen Entropierate verknüpft; vgl.\ \cite{CoverThomas2006}.
\end{proof}

%% ============================================================
\section{Literatur}
%% ============================================================

\begin{thebibliography}{99}

\bibitem{Shannon1948}
C.~E. Shannon,
\emph{A Mathematical Theory of Communication},
Bell System Technical Journal \textbf{27} (1948), 379--423 und 623--656.

\bibitem{Shannon1959}
C.~E. Shannon,
\emph{Coding Theorems for a Discrete Source with a Fidelity Criterion},
IRE National Convention Record \textbf{4} (1959), 142--163.

\bibitem{CoverThomas2006}
T.~M. Cover und J.~A. Thomas,
\emph{Elements of Information Theory}, 2.\ Aufl.,
Wiley-Interscience, Hoboken, NJ, 2006.

\bibitem{Gallager1963}
R.~G. Gallager,
\emph{Low-Density Parity-Check Codes},
MIT Press, Cambridge, MA, 1963.

\bibitem{Arikan2009}
E.~Ar\i{}kan,
\emph{Channel Polarization: A Method for Constructing Capacity-Achieving Codes
for Symmetric Binary-Input Memoryless Channels},
IEEE Transactions on Information Theory \textbf{55} (2009), Nr.~7, 3051--3073.

\bibitem{Huffman1952}
D.~A. Huffman,
\emph{A Method for the Construction of Minimum-Redundancy Codes},
Proceedings of the IRE \textbf{40} (1952), Nr.~9, 1098--1101.

\bibitem{Kolmogorov1965}
A.~N. Kolmogorow,
\emph{Drei Ansätze zur mengentheoretischen Definition des Informationsbegriffs},
Probleme der Informationsübertragung \textbf{1} (1965), Nr.~1, 1--7.

\bibitem{RichardsonUrbanke2001}
T.~J. Richardson und R.~L. Urbanke,
\emph{The Capacity of Low-Density Parity-Check Codes under Message-Passing
Decoding},
IEEE Transactions on Information Theory \textbf{47} (2001), Nr.~2, 599--618.

\bibitem{Berrou1993}
C.~Berrou, A.~Glavieux und P.~Thitimajshima,
\emph{Near Shannon Limit Error-Correcting Coding and Decoding: Turbo-Codes},
Proc.\ IEEE ICC 1993, 1064--1070.

\bibitem{Singleton1964}
R.~C. Singleton,
\emph{Maximum Distance $q$-nary Codes},
IEEE Transactions on Information Theory \textbf{10} (1964), Nr.~2, 116--118.

\end{thebibliography}

\end{document}
